
\documentclass[11pt,letterpaper,titlepage]{article}
\usepackage[margin=1in]{geometry}
\usepackage[T1]{fontenc}
\usepackage{lmodern}
\usepackage{microtype}
\usepackage{booktabs}
\usepackage[hidelinks]{hyperref}
%\usepackage{graphicx} % Required for inserting images
\usepackage{hyperref}
\usepackage{amssymb}
\usepackage{amsmath}
\usepackage{amsthm}
\usepackage{mathrsfs}
\usepackage{comment}

\def\le{\leqslant}
\def\ge{\geqslant}

\def\Z{\mathbb{Z}}
\def\N{\mathbb{N}}
\def\R{\mathbb{R}}
\def\C{\mathcal{C}}
\def\Q{\mathbb{Q}}
\def\F{\mathbb{F}}
\def\I{\mathcal{I}}
\def\B{\mathcal{B}}

\DeclareMathOperator{\cl}{cl}
\newcommand{\si}{\operatorname{si}}

%\usepackage[margin=1in]{geometry}

\title{CS3263 Report}

\author{
Deng Tianle \\
A0282801J \\
\texttt{e1156929@u.nus.edu}
\and
Liu Lekang \\
A0281207L \\
\texttt{e1155335@u.nus.edu}
\and
Zhao Kejun \\
A0282631H \\
\texttt{e1156759@u.nus.edu}
}

\date{13 November 2026}

\begin{document}

\newtheorem{theorem}{Theorem}[section]
\newtheorem{proposition}[theorem]{Proposition}
\newtheorem{lemma}[theorem]{Lemma}
\newtheorem{corollary}[theorem]{Corollary}

%\theoremstyle{definition}
\newtheorem{definition}{Definition}[section]

\theoremstyle{remark}
\newtheorem*{remark}{Remark}
\newtheorem*{example}{Example}

\maketitle
\section{Introduction}

A matroid records which sets of elements are independent. When those dependencies arise from the columns of a matrix over \(\F\), finite-field linear algebra gives a concrete way to work with the matroid. An \emph{excluded minor} is a minimal obstruction under deletion and contraction to obtaining such a representation. Brettell enumerated the excluded minors for \(\F\)-representability on ten elements \cite{brettell}. This project studies the next size at rank four. The exhaustive computation in \texttt{matroid1} gives 35 isomorphism classes and provides a catalogue and verification records for the result \cite{repo,verification}.

The project also asks how an LLM should contribute to a mathematical search. We compare four workflows, from direct reasoning to translation followed by exact solving. The first experiment asks the workflows to classify supplied matroids. The second asks them to find every valid extension of a supplied parent. Finding valid candidates and establishing that the list is complete require different evidence, and the two experiments measure these requirements separately.

\section{Background and problem statement}

A \emph{matroid} \(M=(E,\mathcal{I})\) consists of a finite ground set \(E\) and a family of independent sets \(\mathcal{I}\). The empty set belongs to \(\mathcal{I}\), subsets of independent sets are independent, and a smaller independent set can be augmented by an element from a larger one. The maximal independent sets are the \emph{bases} of \(M\). They have a common size, the \emph{rank} \(r(M)\), and determine \(\mathcal{I}\). Two matroids are isomorphic if a bijection between their ground sets preserves the bases. The rank of a subset \(S\) is the size of its largest independent subset. Its \emph{closure} contains every element whose addition leaves that rank unchanged. A \emph{hyperplane} is a maximal subset of rank \(r(M)-1\) \cite{oxley}.

Write \(\F=\mathbb{Z}/5\mathbb{Z}\). A rank-\(r\) matroid is \emph{quinary} if there is an \(r\times |E|\) matrix \(V\) over \(\F\), with columns labelled by \(E\), such that for every \(r\)-subset \(S\subseteq E\),
\[
 S\in\B(M)\quad\Longleftrightarrow\quad\det V_S\ne 0.
\]
Deletion \(M\setminus a\) keeps the independent sets avoiding \(a\). For a represented matroid, this removes column \(a\). Contraction \(M/a\) is realized by projecting the other columns into the quotient by the span of column \(a\). A minor is the result of a sequence of these operations. Quinary representability is preserved under both operations, so checking all single-element deletions and contractions suffices to check every proper minor \cite{oxley}.

\paragraph{Problem.}
Enumerate, up to matroid isomorphism, all matroids \(M\) with \(|E(M)|=11\) and \(r(M)=4\) such that \(M\) is nonquinary and both \(M\setminus a\) and \(M/a\) are quinary for every \(a\in E(M)\). The output is a list of exact basis families, together with evidence of excluded-minor status, pairwise distinctness, and completeness of the list.

\section{Method}
\label{sec:method}

We view each target matroid \(M\) as a single-element extension of a ten-element quinary \emph{parent} \(N=M\setminus e\). Following the classification paper \cite{classification-paper}, we use three stages.
\begin{enumerate}
\item Generate every possible parent up to matroid isomorphism.
\item For each parent, enumerate all nonquinary extensions whose single-element deletions are quinary.
\item Check every single-element contraction and remove isomorphic repeats.
\end{enumerate}
The second stage is a finite constraint satisfaction problem: each deletion supplies an allowed table, and the tables must agree on shared basis decisions.

\subsection{Step 1: generating deletion parents}

\begin{lemma}[Possible deletion parents]
\label{lem:method-parents}
If \(M\) is a rank-four, eleven-element excluded minor for quinary representability, then for every \(e\in E(M)\), the deletion \(N=M\setminus e\) is a connected simple rank-four quinary matroid on ten elements.
\end{lemma}
\begin{proof}
An excluded minor for representability over a field on at least four elements is 3-connected. A disconnected matroid decomposes as a direct sum, and a connected matroid with a 2-separation decomposes as a 2-sum. The summands are isomorphic to proper minors and hence are quinary. Direct sums preserve representability by taking block diagonal matrices. For a 2-sum, choose representations in which the common element is a unit vector, identify its one-dimensional span in the two vector spaces, and delete its column. Closure under these sums would make \(M\) quinary, contradicting the excluded-minor condition \cite{oxley}.

Since \(M\) is 3-connected, it has no loops, parallel pairs, or coloops. Therefore \(N\) is simple and has rank four, and it is quinary by minor-minimality. If \(N\) were a direct sum on nonempty sets \(A,B\), name them so that \(|A|\geq 2\). Rank is unchanged by deletion on subsets of \(N\), so \(r_M(A)+r_M(B)=r(N)=r(M)\). Adding \(e\) increases the rank of \(B\) by at most one, giving
\[
 r_M(A)+r_M(B\cup\{e\})-r(M)\leq 1.
\]
Both sides of the partition \((A,B\cup\{e\})\) have at least two elements. This contradicts 3-connectivity of \(M\).
\end{proof}

We search all connected simple rank-four quinary parents. Such a parent has no coloop, since a coloop would give a one-element direct summand. In particular, every further deletion of a parent still has rank four.

\subsubsection{Completeness of represented generation}

A simple rank-\(r\) representation uses distinct points of \(\mathrm{PG}(r-1,5)\), the \((5^r-1)/4\) one-dimensional subspaces of \(\F^r\). We choose one vector per point by scaling its first nonzero entry to 1. We start at \(I_r\), append every unused projective point to each retained matrix, and retain one matrix per \emph{field-isomorphism} class at each size. Field isomorphism allows invertible row operations, nonzero column scalings, and column permutations. SageMath tests this equivalence.

\begin{proposition}[Coverage at each size]
\label{prop:method-generation}
At each generated size \(k\geq r\), the procedure retains a representative of every field-isomorphism class of simple rank-\(r\) quinary representations on \(k\) elements.
\end{proposition}
\begin{proof}
Take any target representation \(W\), move a basis to the first \(r\) columns, and reduce it to \(I_r\). Let \(W_k\) be the first \(k\) columns of this matrix. The seed covers \(W_r\). If a retained matrix \(V_k\) is field-isomorphic to \(W_k\), write \(W_{k+1}=[W_k\mid w]\) and let \(G\) be the invertible row transformation used in that equivalence. Extending the old column operations to fix the new column shows that \([V_k\mid Gw]\) is field-isomorphic to \(W_{k+1}\). The direction of \(Gw\) is unused in \(V_k\), since otherwise \(w\) would be parallel to an old column of \(W_k\). The procedure therefore tries this point and retains its class. Induction reaches \(W\).
\end{proof}

For rank four, there are 156 projective points. At each size we can convert the matrices to basis families by testing all four-column determinants and deduplicate by exact matroid isomorphism. The abstract catalogue stores one bit per four-subset in lexicographic order, so a ten-element parent uses \(\binom{10}{4}=210\) bits. The represented catalogue is retained separately because a single abstract matroid can have several field-inequivalent representations with different possible extensions. We keep one labelled representative of each connected abstract class at size ten. We also retain the size-nine represented catalogue for the deletion tables and generate the simple rank-three catalogue through size ten for contraction tests.

\subsection{Step 2: recovering extensions from compatible deletions}

Fix a parent \(N\) on \(P=\{0,\ldots,9\}\) and give the new element the label \(e=10\). Let \(\mathcal{D}(N)\) be the set of simple rank-four matroids \(M\) on \(P\cup\{e\}\) such that \(M\setminus e=N\), \(M\) is nonquinary, and every single-element deletion of \(M\) is quinary. We call these the \emph{deletion-minimal nonquinary extensions} of \(N\).

The bases avoiding \(e\) are fixed by \(N\). The 120 Boolean variables \(x_T\), indexed by triples \(T\in\binom{P}{3}\), specify the remaining basis decisions. The proposed basis family is
\[
 \B_N(x)=\B(N)\cup\bigl\{T\cup\{e\}:T\in\textstyle\binom{P}{3},\ x_T=1\bigr\}.
\]
For \(Q=N\) or \(Q=N\setminus f\), let \(R(Q)\) contain the bit rows arising from all simple rank-four quinary extensions of \(Q\) by \(e\). Let \(\pi_f(x)\) retain the coordinates for triples avoiding \(f\). For instance, deletions of 3 and 4 must agree on \(x_{\{0,1,2\}}\), the status of \(\{0,1,2,10\}\). This bit is visible only in deletions outside that four-set. Requiring the proposed deletions to occur in their tables and excluding the full-parent rows gives
\[
 S(N)=\bigl\{x\in\{0,1\}^{\binom{P}{3}}:
 \pi_f(x)\in R(N\setminus f)\text{ for every }f\in P\bigr\}
       \setminus R(N).
\]

\begin{lemma}[Deletion compatibility]
\label{lem:method-compatibility}
Let \(|E|>2r\) and let \(\mathcal{B}\) be a nonempty family of \(r\)-subsets of \(E\). If for every \(a\in E\) the members avoiding \(a\) form the bases of a rank-\(r\) matroid on \(E\setminus\{a\}\), then \(\mathcal{B}\) is the basis family of a rank-\(r\) matroid on \(E\).
\end{lemma}
\begin{proof}
Take \(U,V\in\mathcal{B}\) and \(u\in U\setminus V\). Since \(|U\cup V|\leq 2r<|E|\), choose \(a\notin U\cup V\). Basis exchange in the specified deletion at \(a\) gives \(v\in V\setminus U\) such that \((U\setminus\{u\})\cup\{v\}\) belongs to that deletion's basis family, and hence to \(\mathcal{B}\). This proves basis exchange for the nonempty family \(\mathcal{B}\).
\end{proof}

\begin{proposition}[Exact recovery]
\label{prop:method-recovery}
The map sending \(x\in S(N)\) to the matroid with bases \(\B_N(x)\) is a bijection onto \(\mathcal{D}(N)\).
\end{proposition}
\begin{proof}
The family \(\B_N(x)\) contains \(\B(N)\). Its restriction avoiding \(e\) gives the bases of \(N\), and its restriction avoiding \(f\in P\) gives the bases of a simple rank-four quinary extension of \(N\setminus f\) by the condition on \(\pi_f(x)\). Since \(11>2\cdot 4\), Lemma~\ref{lem:method-compatibility} makes \(\B_N(x)\) a matroid basis family. The deletions have rank four and are simple, so every pair of elements is independent in a deletion containing it, and hence in the recovered matroid. The matroid is nonquinary because its row is outside \(R(N)\).

Conversely, if \(M\in\mathcal{D}(N)\), then each \(M\setminus f\) for \(f\in P\) is simple and quinary. It has rank four because it contains \(N\setminus f\), which has rank four. Its extra basis decisions therefore lie in \(R(N\setminus f)\). Since \(M\) is nonquinary, its full row lies outside \(R(N)\). This row determines \(M\) uniquely because \(\B(N)\) and the bits \(x_T\) together specify all bases.
\end{proof}

\subsubsection{Labelled representations}

To build \(R(Q)\) exhaustively, we first construct \(L(Q)\), one matrix per \emph{labelled projective equivalence} class representing \(Q\). This equivalence allows invertible row operations and nonzero column scalings with the element labels fixed. For each field-isomorphism class representing a matroid isomorphic to \(Q\), we align one catalogue matrix with \(Q\) using a matroid isomorphism \(\alpha\).

To see why the aligned matrices suffice after relabelling, take any labelled representation of \(Q\). Proposition~\ref{prop:method-generation} provides a field-isomorphic catalogue matrix, and the column bijection \(\beta\) in that field isomorphism preserves the bases. Both \(\alpha\) and \(\beta\) are matroid isomorphisms onto \(Q\), so \(\beta\circ\alpha^{-1}\) is an automorphism of \(Q\). Applying this automorphism to the aligned matrix reaches the target class up to row operations and column scalings.

We compute the automorphisms from a bipartite incidence graph, with one colour for element vertices and another for basis vertices. An edge records membership of an element in a basis. A colour-preserving automorphism permutes the elements and preserves the basis family, and every matroid automorphism extends to the basis vertices. SageMath supplies group generators. Repeatedly applying them to the aligned matrices and retaining each new normalized matrix exhausts the orbits. Relabelling respects projective equivalence, so the normalized matrix suffices for continuing the orbit calculation.

For normalization, fix an ordered basis \(B=(b_1,\ldots,b_r)\), reduce its columns to \(I_r\), and write the remaining block as \(Z\). The determinant for a single basis exchange gives
\[
 Z_{ij}\ne0\quad\Longleftrightarrow\quad
 (B\setminus\{b_i\})\cup\{j\}\in\B(Q),\qquad j\in E(Q)\setminus B.
\]

\begin{lemma}[Normalization by a support forest]
\label{lem:method-normalization}
With the ordered basis fixed, scaling all entries on a fixed spanning forest of the support graph to 1 gives a unique matrix in each labelled projective equivalence class.
\end{lemma}
\begin{proof}
In the support graph, row vertex \(i\) is adjacent to column vertex \(j\) when \(Z_{ij}\ne 0\). Choose a root multiplier of 1 in each tree. Traversing the forest, each edge determines the multiplier at its unvisited endpoint that makes the entry 1. Inverse scalings on the basis columns preserve \(I_r\).

For uniqueness, a row transformation between two equivalent matrices in identity form must preserve each basis-column direction and is therefore diagonal. If both matrices have forest entries 1, the product of the row and column multipliers on each forest edge must equal 1. It follows along paths that all row multipliers in each tree are equal and all column multipliers are their inverses. Every support edge has both endpoints in the same tree, so every nonzero entry is unchanged. Isolated vertices affect no entry of \(Z\).
\end{proof}

\subsubsection{Extension tables and hyperplane compression}

For every \(V\in L(Q)\), append each unused point \(v\in\mathrm{PG}(3,5)\) as column \(e\). The resulting matrix represents a simple quinary extension, with row
\[
 z(V,v)_T=\mathbf{1}\{\det[V_T\mid v]\ne0\},
 \qquad T\in\binom{E(Q)}{3}.
\]
Every resulting row lies in \(R(Q)\). Conversely, take a row in \(R(Q)\) and a representation \([W\mid w]\) of its extension, with \(w\) labelled \(e\). Since \(W\) represents \(Q\), some \(V\in L(Q)\) is equivalent to \(W\) with labels fixed. Apply the row transformation to all columns and the old-column scalings to \(W\). Normalize the transformed new column \(v\) by scaling it. The direction of \(v\) is unused because the extension is simple, so \([V\mid v]\) is tried and yields the original row. Thus the distinct rows are exactly \(R(Q)\).

To batch these tests, expand each determinant along \(v\) as \(c_T^{\mathsf T}v\). Stacking the cofactor rows in \(C\) and all 156 projective point columns in \(D\), the product \(CD\) over \(\F\) evaluates every triple--point test for that matrix at once. We omit used points and convert nonzero entries to 1.

In any rank-four extension \(M\) of \(N\), a dependent triple \(T\) has bit 0 because \(r_M(T\cup\{e\})\leq 3\). If \(T\) is independent, its closure \(H=\cl_N(T)\) is a hyperplane. Since \(N\) is a restriction of \(M\), we have \(T\subseteq H\subseteq\cl_M(T)\), so \(\cl_M(H)=\cl_M(T)\). Thus \(T\cup\{e\}\) is a basis exactly when \(e\notin\cl_M(H)\). All triples spanning \(H\) share this decision, which we denote by \(y_H\).

Let \(\mathcal{H}(N)\) denote the hyperplanes of \(N\). The coordinates visible after deleting \(f\) are
\[
 C_f=\left\{\operatorname{cl}_N(T):
 T\in\binom{P\setminus\{f\}}{3},\ r_N(T)=3\right\}.
\]
For each surviving independent triple, \(\cl_{N\setminus f}(T)=\cl_N(T)\setminus\{f\}\). Thus triples grouped by a coordinate in \(C_f\) share a closure in \(N\setminus f\) and have equal bits in every row of \(R(N\setminus f)\). Dependent triples stay dependent under deletion. We therefore compress \(R(N\setminus f)\) to a table \(A_f\) on \(C_f\) and \(R(N)\) to a table \(A\) on \(\mathcal{H}(N)\). Define
\[
 \widehat S(N)=
 \bigl\{y\in\{0,1\}^{\mathcal{H}(N)}:
 y|_{C_f}\in A_f\text{ for every }f\in P\bigr\}\setminus A.
\]
Expanding a row \(y\) copies \(y_H\) to every independent triple spanning \(H\) and puts zero on dependent triples. This reverses the compression in each table. Hence \(y|_{C_f}\in A_f\) exactly when the expanded row satisfies \(\pi_f(x)\in R(N\setminus f)\), and \(y\in A\) exactly when \(x\in R(N)\). Conversely, each \(x\in S(N)\) gives a matroid by Proposition~\ref{prop:method-recovery}, so its bits are constant on each hyperplane group and compress to a unique \(y\). This gives a bijection between \(\widehat S(N)\) and \(S(N)\).

\subsubsection{Enumeration by natural joins}

A natural join combines two table rows exactly when they agree on every shared coordinate. For a nonempty set \(I\subseteq P\) of deletions already joined, let \(C_I=\bigcup_{f\in I}C_f\). The accumulated relation is
\[
 J_I=\{y\in\{0,1\}^{C_I}:y|_{C_f}\in A_f\text{ for all }f\in I\}.
\]
For a single deletion this is its table. Joining \(J_I\) with the next table enforces the previous conditions on \(C_I\) and the new condition on that table's columns, proving the formula inductively. Every hyperplane \(H\) appears in at least one table because deleting an \(f\notin H\) leaves a spanning triple of \(H\) intact. The joined rows therefore assign every variable, and removing the full-parent rows gives
\[
 \widehat S(N)=(A_0\bowtie A_1\bowtie\cdots\bowtie A_9)\setminus A.
\]
The implementation groups rows by their bits on the shared columns and chooses the next table by the exact size of the resulting join. Once all columns have appeared, rows in \(A\) can be removed immediately because further joins only filter full assignments. Every connected parent in the eleven-element search was exhausted by this method.

\subsubsection{SAT encoding and enumeration}

The same tables admit a propositional SAT encoding with one Boolean variable \(y_H\) per hyperplane. To construct a binary decision diagram (BDD) for \(A_f\), split its rows on successive columns in a fixed order. An empty branch rejects, and a branch that reaches the end with an allowed row accepts. Identical subdiagrams can be shared. For example, the table \(\{00,01,11\}\) on \((y_1,y_2)\) accepts immediately when \(y_1=0\). When \(y_1=1\), it accepts exactly when \(y_2=1\).

At a node testing \(y_H\), let \(v\) be the Boolean value of the subdiagram, and let \(\ell,h\) be the values of its 0- and 1-branches. The node equivalence and its four clauses in conjunctive normal form are
\[
\begin{gathered}
 v\ \longleftrightarrow\ ((\neg y_H\land\ell)\lor(y_H\land h)),\\[0.6em]
\begin{aligned}
 &(\neg v\lor y_H\lor\ell),
 &&(\neg v\lor\neg y_H\lor h),\\
 &(v\lor y_H\lor\neg\ell),
 &&(v\lor\neg y_H\lor\neg h).
\end{aligned}
\end{gathered}
\]
For a fixed \(y_H\), the clauses force \(v\) to equal the value of the selected child. Starting at the terminals and working up, each node is therefore forced to its evaluation on the hyperplane row. Setting each root to true enforces every deletion constraint. To rule out a quinary full extension, for each row \(a\in A\) we add the blocking clause
\[
 \bigvee_{H\in\mathcal{H}(N)}(y_H\ne a_H).
\]
Here the literal is \(y_H\) when \(a_H=0\), and \(\neg y_H\) when \(a_H=1\). The satisfying assignments projected onto the \(y_H\) variables are exactly \(\widehat S(N)\). We record each resulting row and add the same kind of clause on the \(y_H\) variables to block it. Enumeration is complete when the solver returns \texttt{UNSAT}. A resource limit returns \texttt{UNKNOWN} and leaves the search unresolved.

\subsection{Step 3: testing contractions and removing duplicates}

For every candidate \(M\in\mathcal{D}(N)\) and every \(a\in E(M)\), the contraction \(M/a\) has rank three because \(a\) is a nonloop. We test its simplification, obtained by deleting loops and retaining one element per parallel class. If \(M/a\) is quinary, its simplification is quinary by closure under deletion. Conversely, a representation of the simplification extends to \(M/a\) by restoring zero columns for loops and scalar copies for parallel elements. Hence it suffices to test the simplification. It has at most ten elements. At size three it is the free rank-three matroid represented by \(I_3\), and at sizes four through ten we compare it by exact isomorphism with the complete simple rank-three catalogue from Step 1.

Candidates passing every contraction test are excluded minors, since Step 2 establishes nonrepresentability and representability of all deletions. SageMath then removes duplicates up to matroid isomorphism, including candidates arising from different parents. An independent check of distinctness compares canonical forms of coloured incidence graphs. For each matroid, the subset vertices represent either all bases or all nonbases of size four, using bases in a tie. Any matroid isomorphism preserves this choice and induces a colour-preserving graph isomorphism. Conversely, the graph determines the chosen subset family up to a relabelling of the elements. A flag records which family was encoded, so the graph and flag together determine the bases.

\subsubsection{Independent certificates}

For each final representative, the positive witnesses are matrices for its eleven deletions and eleven contractions, with columns in the order of the remaining labels. Modular Gaussian elimination checks the basis status of every subset of the appropriate size against the required minor. SageMath performs a second matrix check.

For nonrepresentability, let \(\Phi_M\) be the SAT formula for a quinary representation of a final matroid \(M\). Fix an ordered basis \(B=(b_1,\ldots,b_r)\) and reduce its columns to \(I_r\), writing the remaining block as \(Z\). As in the labelled normalization, a position \(Z_{ij}\) is nonzero exactly when \((B\setminus\{b_i\})\cup\{j\}\) is a basis. Lemma~\ref{lem:method-normalization} allows the forest entries to be fixed to 1, leaving the other nonzero entries with domain \(\{1,2,3,4\}\).

For each \(r\)-subset \(S\subseteq E(M)\), let \(I(S)=\{i:b_i\notin S\}\) and \(J(S)=S\setminus B\). Expanding along the identity columns reduces the determinant on \(S\), up to sign, to the square minor on the missing rows and extra columns. Thus we impose
\[
 \det Z[I(S),J(S)]\ne 0
 \quad\Longleftrightarrow\quad S\in\B(M),
\]
with the empty determinant equal to 1. The encoder computes these determinants by cached Laplace expansion over \(\F\).

Each field expression \(u\) with domain \(D_u\) has Boolean indicators \(p_{u,a}\) for \(a\in D_u\), with exactly one true. For an addition \(w=u+v\), the formula includes
\[
 \neg p_{u,a}\ \lor\ \neg p_{v,b}\ \lor\ p_{w,a+b}
\]
for every allowed pair \(a,b\), with the sum taken modulo five. Multiplication uses the analogous clause with \(ab\). Since the output has exactly one value, the selected inputs force the correct output. Constants can be substituted directly. Induction over the arithmetic operations therefore forces the determinant values, and the formula constrains their indicators according to the specified basis status.

Any representation of \(M\) can be normalized as above and gives a satisfying assignment by evaluating the field expressions. Conversely, a satisfying assignment gives a full-rank matrix in identity form whose nonsingular \(r\)-column submatrices are exactly the bases of \(M\). Hence \(\Phi_M\) is satisfiable exactly when \(M\) is quinary. The saved DRAT refutation for each representative is checked against its saved formula by DRAT-trim \cite{drat}. Together with the 22 minor matrices, this certifies excluded-minor status directly from the representative's basis family.

\subsection{Search coverage and classification}

At size ten, generation retained 64,508 field-isomorphism classes, which gave 58,756 abstract classes. Aligning matrices and closing their automorphism orbits gave 71,547 labelled projective classes across that full catalogue. The merge checked that the four search ranges had matching input and source fingerprints and covered the catalogue without gaps or overlaps. Of the 58,756 parents, 84 were disconnected and every remaining problem finished by natural join with no unresolved cases. The 659 pairs in Table~\ref{tab:enumeration} count labelled extensions together with their parents, before the contraction filter \cite{classification-paper,verification}.

\begin{table}[!htbp]
\centering
\caption{Stages of the exhaustive rank-four search.}
\label{tab:enumeration}
\begin{tabular}{@{}lr@{}}
\toprule
Stage & Count \\
\midrule
Simple quinary ten-element parent classes & 58,756 \\
Connected parent classes searched to exhaustion & 58,672 \\
Labelled pairs \((N,M)\) with \(M\in\mathcal{D}(N)\) & 659 \\
Excluded-minor classes after contractions and deduplication & 35 \\
\bottomrule
\end{tabular}
\end{table}

\begin{theorem}
There are exactly 35 isomorphism classes of rank-four, eleven-element excluded minors for quinary representability.
\end{theorem}
\begin{proof}
For the upper bound, take any excluded minor of the required rank and size. By Lemma~\ref{lem:method-parents} and Proposition~\ref{prop:method-generation}, one of its deletions is isomorphic to a searched parent \(N\). Relabel the deletion as \(N\) and the removed element as \(e\). The resulting matroid belongs to \(\mathcal D(N)\), so Proposition~\ref{prop:method-recovery}, the complete extension tables, and the exhaustive joins ensure it is recovered. All its contractions are quinary, so its isomorphism class survives Step 3.

For the lower bound, the 35 checked DRAT refutations establish nonquinary status, and the \(35\times 22=770\) checked matrices establish representability of all immediate minors. The independent incidence-graph check confirms that these representatives are pairwise nonisomorphic \cite{verification}.
\end{proof}

As calibration, the implementation reproduced Brettell's rank-four catalogues by exact matroid isomorphism, with 5, 92, 219, and 130 classes at sizes seven through ten, respectively \cite{brettell,verification}. At sizes seven and eight, the condition \(|E|>2r\) fails, so the searches impose basis exchange explicitly. SAT and natural join also agreed on the solutions from the shared tables for every connected parent at target size nine and on all 5,991 labelled deletion-minimal candidates at target size ten \cite{classification-paper,verification}.

The repository's \path{results/r4-n11/} directory contains the merged search record (\path{summary.json}), the exact basis families (\path{excluded.jsonl}), and the 770 witness matrices (\path{minors/minor-witnesses.jsonl}). Each case under \path{nonrepresentation/} includes the DIMACS formula, DRAT refutation, and checker output. The generation, search, and replay commands are documented in \path{docs/REPRODUCING.md} \cite{repo}.


\section{Experiments and results}

We evaluate how an LLM should contribute to the classification task through two experiments. The first tests classification of supplied matroids; the second tests the more demanding task of finding every excluded-minor extension of a supplied parent. Together they distinguish answer correctness, computational efficiency, and evidence of completeness.

\subsection{Workflows and shared setup}

We compare four workflows with different divisions of responsibility between the frontier LLM model and the tools and solvers:
\begin{description}
\item[Direct reasoning.] The model returns its answer and justification in one call, without tools.
\item[Exact queries.] The model requests ranks, determinants modulo five, deletions, or contractions, then interprets these calculations and constructs the final answer.
\item[Model-directed SAT.] The model chooses representability tests in classification, or candidate checks and constrained extension searches in enumeration. It uses the solver replies to decide the final label or solution set.
\item[Translation and exact solving.] The model translates the supplied hyperplanes into structured data. After checking that the translation preserves the full labelled basis family, the exact backend performs the classification or exhaustive search and returns the answer.
\end{description}

Both experiments use \texttt{gpt-6.1-sol} at medium reasoning effort, with one trial per instance and workflow. All workflows receive the same complete hyperplane descriptions, while answer keys remain with the grader. Built-in model tools are disabled; only the designated interfaces are available. The benchmark cases are separate from the pilots. Query budgets, case-generation details, and supplementary results are provided in Appendix~\ref{app:experiments}.

\subsection{Experiment 1: classifying supplied matroids}

\paragraph{Design.}
We use 30 rank-four matroids on eleven elements, with ten in each class: quinary, excluded minor, and nonquinary but nonminimal. Quinary cases are generated from matrices, excluded minors are drawn from the certified catalogue, and nonminimal cases extend known excluded minors so that deletion recovers an obstruction. All cases are relabelled and their labels checked by the exact solver. We measure classification accuracy and median elapsed time. For the ten quinary cases, we also check matrices supplied by the direct and exact-query workflows against every basis and nonbasis. This separates a correct label from a verifiable representation.

\begin{table}[!htbp]
\centering
\caption{Classification results on 30 supplied matroids.}
\label{tab:classification}
\begin{tabular}{@{}lrr@{}}
\toprule
Workflow & Correct labels & Median time (s) \\
\midrule
Direct reasoning & \(30/30\) & 50.8 \\
Exact queries & \(30/30\) & 80.2 \\
Model-directed SAT & \(30/30\) & 16.3 \\
Translation and exact solving & \(30/30\) & 12.2 \\
\bottomrule
\end{tabular}
\end{table}

\paragraph{Results.}
All four workflows classify every case correctly (Table~\ref{tab:classification}). Translation and exact solving has the shortest median time, about one quarter of direct reasoning, and all 30 translations preserve the input basis family. All ten matrices from direct reasoning pass the determinant checks, as do nine from exact queries. The remaining exact-query response gives the correct label but includes a placeholder instead of a matrix.

\paragraph{Discussion.}
Exact queries increase elapsed time without improving label accuracy. Computing a deletion or contraction still leaves the model to decide its representability; SAT queries resolve that predicate directly. The extra interaction also contributes to time. Thus the useful tool interface is the one that addresses the remaining mathematical decision, rather than simply supplying intermediate calculations. The missing matrix further shows why accuracy should be accompanied by a certificate check.

The perfect scores require qualification. The case constructions induce different basis-count distributions: a leave-one-out rule assigning the nearest class mean from basis count alone scores \(24/30\). All workflows solve the six cases it misses, but the benchmark still supplies a strong structural cue. Relabelling preserves that cue. These results demonstrate success on the selected cases; a stronger generalization test would match classes on basis counts, diversify their constructions, and repeat trials.

\subsection{Experiment 2: searching for all excluded-minor extensions}
\label{sec:search}

\paragraph{Design.}
Each trial supplies a connected, simple, rank-four quinary parent on ten elements and asks for all excluded-minor extensions by element 10. The 24 pairwise nonisomorphic parents comprise three groups of eight: parents with valid solutions, parents whose deletion-compatible nonquinary candidates fail contraction checks, and parents with no such candidates. This design tests recovery and two distinct reasons for an empty answer. Five positive parents have multiple solutions. Exhaustive joins, SAT enumeration, contraction checks, and comparison with the certified catalogue establish 16 reference solutions up to isomorphism fixing element 10.

We measure candidate precision, solution recall, and \emph{exact-set match}, which requires a complete, valid answer for a parent. We separately measure \emph{verified exhaustion}: an accepted computational trace must cover all extensions and account for every valid output. Matching the reference set and establishing completeness are different requirements. Median time covers the workflow's search and verification. Prepared compatibility tables are available to the SAT and translation workflows; their median construction time of 0.3 seconds is recorded separately.

\begin{table}[!htbp]
\centering
\caption{Search results on 24 parents. Exhaustion counts accepted computational traces; free-text completeness arguments are not scored by this measure.}
\label{tab:search-main}
\begin{tabular}{@{}lrrrrr@{}}
\toprule
Workflow & Precision & Recall & Exact sets & Exhausted & Median (s) \\
\midrule
Direct reasoning & \(16/16\) & \(16/16\) & \(24/24\) & \(0/24\) & 131.1 \\
Exact queries & \(16/16\) & \(16/16\) & \(24/24\) & \(0/24\) & 144.8 \\
Model-directed SAT & \(15/16\) & \(15/16\) & \(23/24\) & \(14/24\) & 79.9 \\
Translation and exact solving & \(16/16\) & \(16/16\) & \(24/24\) & \(24/24\) & 9.9 \\
\bottomrule
\end{tabular}
\end{table}

\paragraph{Results.}
Direct reasoning, exact queries, and translation recover all 16 solutions and match every reference set (Table~\ref{tab:search-main}). Translation is fastest on every parent and provides accepted exhaustion evidence in all 24 trials. Model-directed SAT supplies correct sets in 23 trials, but only 14 have accepted exhaustion traces; nine correct answers leave coverage gaps. In its remaining trial, the solver finds both required extensions, but the final response omits a triple from one candidate, making the submitted basis family invalid. All workflows correctly return empty sets on the 16 negative parents.

\paragraph{Discussion.}
The distinction between finding solutions and proving completeness is decisive. Direct reasoning succeeds at recovering the answers, but its free-text arguments are not accepted computational traces. Its \(0/24\) exhaustion score therefore does not establish that those arguments are mathematically wrong. Model-directed SAT provides exact local answers, yet still depends on the model to cover every branch and preserve checked candidates. The transcription failure demonstrates that solver correctness need not survive the final response. Returning identifiers for checked basis families would address this failure mode; that interface improvement remains to be evaluated.

Translation transfers enumeration, contraction checks, and output construction to the backend. Its 9.9-second median is about thirteen times lower than direct reasoning's, while also providing complete traces. The exact method exploits the problem's structure: these parents require 28--48 hyperplane variables, but the largest intermediate join contains only 1,166 rows. Compatibility constraints sharply restrict the represented search space, although this is not a measured speedup over an uncompressed baseline. The negative groups also confirm the need for both filters: each contraction near miss leaves one deletion-compatible candidate that must still be rejected.

The balanced sample makes these behaviours observable but does not reflect their frequency in the full search. Only 224 of 58,672 connected parents, about \(0.38\%\), admit excluded-minor extensions. An always-empty answer would match about \(99.62\%\) of that catalogue while discovering nothing. Consequently, solution recall and exhaustion evidence are essential alongside exact-set accuracy. The benchmark evaluates selected search behaviours; the global classification rests on the separate exhaustive computation.

\subsection{Overall Discussion}

Across both experiments, validated translation followed by exact solving gives the fastest answers and the strongest computational evidence. This supports assigning formal verification and exhaustive enumeration to exact tools for fully specified instances. It does not demonstrate a need for an LLM to parse the present regular format: a deterministic parser is an appropriate unmeasured baseline. Ambiguous or incomplete descriptions would require specification checks before a solver's certificate could support the intended claim.

These are comparisons of complete workflows, whose tool capabilities and model-call counts differ, rather than controlled component ablations. Timings exclude catalogue construction and the separately recorded search-table preparation. With one trial per instance, one model, and constructed cases, the results do not establish statistical stability or performance on larger or less structured inputs. The supported conclusion concerns the division of work in this setting, rather than a general superiority of tool use over direct reasoning.




\begin{comment}
\section{Experimental study}

\subsection{Instance descriptions and validation}
\label{sec:inputs}

The input must specify a single labelled matroid, rather than a few dependencies that admit several completions. For the two experiments below, each prompt gives the ground set, rank, and all hyperplanes of size at least the rank. Writing this hyperplane list as \(\mathcal H\), the complete basis family is
\[
 \B(M)=\{S\subseteq E: |S|=r(M),\ S\nsubseteq H\text{ for every }H\in\mathcal H\}.
\]
Indeed, every dependent set of size \(r(M)\) lies in a hyperplane, and that hyperplane has at least \(r(M)\) elements. This compact description therefore fixes all \(\binom{11}{4}=330\) basis-membership decisions in a classification instance, or all \(\binom{10}{4}=210\) in a search parent. Checking a translation means comparing these entire labelled families, not merely checking that the output satisfies the matroid axioms.

A complementary diagnostic collection stores the complete bases explicitly. It contains 24 simple rank-four, eleven-element matroids, with eight quinary, eight excluded-minor, and eight nonquinary nonminimal cases. Six are designated for development and 18 for testing, with each split balanced across the three classes. The collection uses explicit matrix constructions, catalogue representatives, and free extensions of nonquinary parents, as in the classification design below. Case identifiers and element permutations conceal the construction labels, while the evaluator retains the provenance and inverse correspondence. All 24 have distinct label-invariant signatures: basis counts, rank distributions of subsets of sizes one through three, and sorted element--basis incidence counts. Unequal signatures prove that no two cases are isomorphic; equality would require a further isomorphism test \cite{diagnostic}.

This collection provides an additional check of the descriptions and answer certificates. Each exported basis list was parsed back and compared with its structured input; all 24 families passed basis exchange and the rank, size, and simplicity checks. Rechecking the labels required eight representation matrices, eight original nonrepresentation proofs with \(8\times22=176\) immediate-minor matrices, and eight nonrepresentable-minor proofs. All checks passed. The determinant formulas were regenerated before replaying DRAT, binding each refutation to its intended matroid. These are checks of a separate diagnostic collection, not additional model trials or part of either benchmark reported below \cite{diagnostic}.

\subsection{Shared protocol}

Both experiments use \texttt{gpt-6.1-sol} at medium reasoning effort through the Codex CLI. Every invocation starts a fresh session. Follow-up prompts repeat the instance and include the earlier queries and tool replies from that trial. The four workflows see the same cases, and the runner disables the model's built-in tools. The model receives only the description and permitted replies. The answer keys are held by the grader. The hyperplane descriptions specify the complete inputs as described in Section~\ref{sec:inputs}. We score one trial per instance and workflow.

\begin{description}
\item[Direct reasoning.] The model reads the instance and returns its answer with a justification in a single call.
\item[Exact queries.] The model can request ranks, determinants modulo five, deletions, or contractions. The tools return these calculations, and the model decides what they imply.
\item[Model-directed SAT.] The model chooses representability checks in the classification experiment. In the search experiment it can test a proposed extension or constrain the status of selected triples. A SAT query on those constraints returns at most one compatible completion, with a contraction verdict. Previously returned completions are blocked.
\item[Translation and exact solving.] The model converts the description to structured hyperplanes. The harness compares the resulting labelled basis family with the original and then hands it to the exact backend. The backend produces the final decision or list.
\end{description}

\paragraph{Answer and evidence requirements.}
For classification, a quinary answer can be certified by a matrix whose maximal determinants match the entire basis family. An excluded-minor answer needs nonrepresentability of the original together with representations of all 22 immediate minors. A nonquinary nonminimal answer can instead be established by one nonrepresentable immediate minor, since representability is minor-closed. A rejected candidate matrix or an unsuccessful bounded search establishes none of these negative claims. For extension search, an empty list is likewise a complete answer only when every possible extension has been ruled out. These distinctions separate label accuracy, valid witnesses, and verified exhaustion in the results below; free-text mathematical arguments remain available for review rather than being treated as mechanically checked certificates.

\subsection{Experiment 1. Classifying supplied matroids}

\paragraph{Design.}
We used 30 rank-four matroids on eleven elements, with ten in each of three classes. The quinary cases came from matrices obtained by extending \(I_4\) with sampled projective points. The excluded minors came from the certified catalogue. For the nonquinary, nonminimal cases, we added an element to a known ten-element excluded minor. The element was added freely to the whole matroid or to a chosen hyperplane, so deleting it recovered the known obstruction. We relabelled the matroids and verified their classification labels with the exact solver. The benchmark cases were separate from the three-case pilot.

For the whole-matroid free extension of a rank-four parent \(N\) by \(e\), the basis family has the explicit form
\[
 \B(N+e)=\B(N)\ \cup\
 \{T\cup\{e\}:T\subseteq E(N),\ |T|=3,\ r_N(T)=3\}.
\]
Deleting \(e\) recovers \(N\), so a nonquinary parent supplies a concrete proper-minor obstruction. The construction preserves rank and, for these simple parents, simplicity. Thus nonminimality need not be signalled by an easily detected loop or parallel pair. The hyperplane extensions in the benchmark retain the same deletion obstruction, but use a different restriction on the new bases.

The two query workflows each had one batch of at most 23 requests, enough for the model-directed SAT workflow to query the original and all 22 immediate minors. Each solver call had a 10-second budget. We measured classification accuracy and median elapsed time per trial. We also checked the matrices supplied by the direct and exact-query models on the ten quinary cases. This witness check measures whether the response provides verifiable evidence in addition to a correct label.

\begin{table}[!htbp]
\centering
% REFERENCE TO SUPPLY: classificationlogs has no bibliography entry in the supplied report.
% This citation is for the 30-case classification experiment, not study/results/b0-linux.
\caption{Classification on 30 held-out matroids \cite{classificationlogs}.}
\label{tab:classification}
\begin{tabular}{@{}lrrr@{}}
\toprule
Workflow & Correct & Median time (s) & Model calls \\
\midrule
Direct reasoning & \(30/30\) & 50.8 & 30 \\
Exact queries & \(30/30\) & 80.2 & 60 \\
Model-directed SAT & \(30/30\) & 16.3 & 60 \\
Translation and exact solving & \(30/30\) & 12.2 & 30 \\
\bottomrule
\end{tabular}
\end{table}

\paragraph{Results and interpretation.}
All four workflows reached the same accuracy (Table~\ref{tab:classification}). Translation and exact solving had the shortest median time, about one quarter of direct reasoning. All 30 translations recovered the original basis family. Model-directed SAT requested all 23 checks on every case, including the representable cases for which a checked matrix for the original already settles the decision. SAT queries resolve representability directly. The exact-query workflow mainly requested deletions and contractions and still had to reason about whether they were representable. Its extra round trip and the remaining proof work help explain why it took longer.

All ten matrices supplied by direct reasoning passed the determinant checks, as did nine of the ten from exact queries. The remaining exact-query response gave the correct label with a placeholder where a matrix should have appeared. The accuracy metric credits this response, while the witness check records the missing evidence.

The construction also gives the classes different basis counts. The ten excluded minors have 233--261 bases, the quinary cases 260--284, and the nonminimal cases 275--293. A leave-one-out rule that predicts the class with the nearest mean basis count scored \(24/30\). Each prediction used the other 29 cases to estimate the three means. All workflows solved the six cases this rule missed. The benchmark therefore offers a strong structural cue, and the matrix audit is a more discriminating measure of the evidence supplied by the model.

\subsection{Experiment 2. Searching for excluded-minor extensions}
\label{sec:search}


\paragraph{Design.}
Each trial gives a connected, simple, rank-four quinary parent \(N\) on ten elements. The task is to return every excluded-minor extension by element \(10\), with isomorphisms required to fix that new element. The response gives the old triples that form bases together with \(10\). This determines each extension because \(N\) fixes the old bases.

The held-out set has 24 pairwise nonisomorphic parents, with eight in each group. The first group comes from deletions of catalogue matroids and has valid solutions. Five of those parents have multiple solutions. The second group comes from deletions of series extensions of a rank-three excluded minor. A series extension here adds an element whose contraction recovers that obstruction. Its deletion-compatible candidates fail the contraction test. The third group comes from replacing a column in a quinary representation of a parent with solutions. For the selected parents in this group, the compatibility search produces no nonquinary candidates. Thus the two groups with empty answer sets test different parts of the search.

We exhaust both the join and the SAT encoding on each parent and compare their labelled outputs. The grader checks contractions and confirms the surviving extensions by matching the parents against deletions of the certified catalogue. This gives 16 solutions up to isomorphism fixing \(10\), from 14 ordinary catalogue classes. The three pilot parents were kept separate. The query workflows could request two batches of up to six checks, with a required first batch. Each trial had a 240-second budget, with a 10-second solver budget per call. All 96 trials finished within budget. The SAT and translation workflows used prepared compatibility tables, whose median construction time was 0.3 seconds per parent and was recorded separately.

\paragraph{Metrics.}
Precision counts valid excluded-minor extensions among submitted candidates. Recall counts the recovered solutions out of the 16 in the reference sets, with duplicates identified up to isomorphism fixing \(10\). An \emph{exact-set match} requires a valid, complete answer for one parent, including the empty set where appropriate. Catalogue coverage counts distinct ordinary isomorphism classes across all parents. \emph{Verified exhaustion} requires a checkable trace that accounts for all extensions and confirms that the final answer includes every valid one. The grader checks complete enumeration traces and SAT branch coverage. It retains free-text arguments for review. We also record model calls, median elapsed time, and median time to the first valid candidate on trials that produced one.

\begin{table}[!htbp]
\centering
% REFERENCE TO SUPPLY: searchlogs has no bibliography entry in the supplied report.
% This citation is for the 24-parent extension-search experiment.
\caption{Search on 24 held-out parents. Precision uses submitted candidates, and recall uses the reference solutions.}
\label{tab:search}
\begin{tabular}{@{}lrrrrr@{}}
\toprule
Workflow & Precision & Recall & Exact sets & Exhausted & Median (s) \\
\midrule
Direct reasoning & \(16/16\) & \(16/16\) & \(24/24\) & \(0/24\) & 131.1 \\
Exact queries & \(16/16\) & \(16/16\) & \(24/24\) & \(0/24\) & 144.8 \\
Model-directed SAT & \(15/16\) & \(15/16\) & \(23/24\) & \(14/24\) & 79.9 \\
Translation and exact solving & \(16/16\) & \(16/16\) & \(24/24\) & \(24/24\) & 9.9 \\
\bottomrule
\end{tabular}
\end{table}

\paragraph{Correctness and coverage.}
All workflows returned the empty set on the 16 parents with no solutions. Among the five parents with multiple solutions, model-directed SAT submitted the full set on four, while the other workflows did so on all five. Direct reasoning, exact queries, and translation recovered all 14 reachable catalogue classes. The submitted SAT answers covered 13. The direct and exact-query completeness arguments were free text, so the automated exhaustion score credits only the two workflows that supplied accepted traces.

On parent \texttt{a48cea5e8b60}, the model queried both basis statuses of \(\{0,1,2\}\) with the new element. Both branches returned \texttt{UNSAT}, exhausting all deletion-compatible nonquinary extensions. The model obtained verified exhaustion on all eight parents with no such candidates, three contraction near misses, and three parents with solutions. Nine other trials returned the correct set and left gaps in the branch trace. On parent \texttt{71b692bcf18e}, the trace covered the search and the solver found both valid extensions. The submitted list for one extension omitted \(\{2,7,9\}\), although it retained other triples spanning the same parent hyperplane. These triples must all have the same basis status with element \(10\), so the submitted basis family is invalid. Returning solver candidate identifiers would preserve the checked basis families through the final response.

\paragraph{Efficiency and scope.}
The 96 trials took about 40 minutes with four workers. Direct, exact-query, SAT, and translation workflows used 24, 48, 55, and 24 model calls respectively. Model-directed SAT reduced median time to a first valid candidate from 144.8 seconds for direct reasoning to 21.3 seconds. Translation and exact solving reached 10.1 seconds and gave the fastest total time on each parent. The parent descriptions already specify every basis, so translation is a short step and the solver can carry the candidate data through enumeration and verification. The speed differences describe the complete workflows, including their different model call counts.

The sample deliberately enriches parents with solutions. Deduplicating all single-element deletions of the 35 certified classes gives 224 such parent classes, about \(0.38\%\) of the 58,672 connected parents. They form one third of our benchmark. The two groups with empty answer sets also use specific constructions to exercise the deletion and contraction filters. These choices make recovery and rejection visible in a small experiment. The reported scores describe the selected cases, while the global result comes from exact enumeration over the full parent catalogue.
\end{comment}
\section{Responsible AI}

\paragraph{Scientific reliability.}
In AI-assisted mathematical research, an incorrect classification can propagate into later proofs and impose substantial checking costs on other researchers. Transparency and the prevention of unintended consequences are therefore central responsible-AI concerns in this domain \cite{responsibleai}. A plausible explanation or a correct benchmark answer is insufficient evidence for a classification theorem. Our claim is supported by the coverage argument, exhaustive searches, checked nonrepresentability refutations, immediate-minor matrices, and isomorphism checks. The distinction matters: certificates establish the status of the listed objects, while completeness also requires coverage of every possible parent and extension. A timeout remains unresolved, and a checked result applies only to the mathematical object actually encoded.

\paragraph{Transparency and human responsibility.}
The AI assistance described in Section~\ref{sec:contributions} makes traceability important: polished mathematical prose can obscure how a result was obtained. The saved prompts, outputs, tool replies, and verification records allow readers to trace claims to computational evidence. Such records support scrutiny but do not explain the model's internal reasoning or guarantee that the implementation matches the intended mathematics. The authors remain responsible for the formulation, completeness argument, attribution of prior work, and claims presented to readers. Neither model confidence nor a solver verdict transfers that responsibility to the software.

\paragraph{Responsible evaluation and use.}
Reporting only accuracy or speed could encourage reliance on answers without adequate justification. Section~4 shows that a correct label can lack its promised witness and that a valid solver result can be corrupted when transcribed. We therefore distinguish answer accuracy, validity of submitted objects, and evidence of exhaustion. Input translations are checked against the original basis family; the observed output error also motivates preserving checked candidates through the final response. The constructed benchmarks support conclusions about these workflows on fully specified matroids, not a general claim that the model can independently conduct reliable mathematical research. Ambiguous descriptions require clarification before certification, and broader use would require additional evaluation.



\section{Repository and supplementary material}

The \href{https://github.com/DengTianle/matroid1}{\texttt{matroid1} repository}
contains the classification code in \path{gf5minors/} and \path{scripts/},
the catalogue and verification records in \path{results/r4-n11/}, and
instructions for reproducing the computation in \path{docs/REPRODUCING.md}.
The detailed proof of the main classification theorem is in
\href{../paper/main.pdf}{\texttt{paper/main.pdf}}.
The slides in \href{../presentation.pdf}{\texttt{presentation.pdf}} provide an
overview of the search and its correctness argument.
The accompanying \href{./appendix.pdf}{\texttt{appendix.pdf}} in
\path{final/} collects the matroid theory definitions and results used in the
report.

\section{Statement of Contributions and AI assistance}
\label{sec:contributions}

\begin{description}
\item[Deng Tianle.] \textit{[]}
\item[Liu Lekang.] \textit{[]}
\item[Zhao Kejun.] \textit{[Designed and run experiments, analyzed results]}
\end{description}

\paragraph{AI assistance.}
We used an LLM with computational tools to develop and carry out the exhaustive search. Codex also assisted with the evaluation harness, analysis of saved results, and drafting and editing this report. These activities were AI-assisted; the authors retain responsibility for the submitted work and its claims.

\clearpage
\begin{thebibliography}{99}

\bibitem{what}
J. Oxley,
\textit{What is a matroid?},
Cubo \textbf{5} (2003), 176--215.

\bibitem{oxley}
J. Oxley,
\textit{Matroid Theory},
2nd ed., Oxford University Press, 2011.

\bibitem{brettell}
N. Brettell,
\textit{The excluded minors for $\mathrm{GF}(5)$-representable matroids on ten elements},
Advances in Applied Mathematics \textbf{166} (2025), 102864.
\url{https://doi.org/10.1016/j.aam.2025.102864}

\bibitem{repo}
Deng Tianle,
\textit{matroid1: computational classification and verification artifacts}.
\url{https://github.com/DengTianle/matroid1}.

\bibitem{classification-paper}
\textit{The rank-four, eleven-element excluded minors for representability over \(\mathbb{F}_5\)}.
Project paper, revised exposition, 5 November 2026, \path{matroid1/paper/main.pdf}.

\bibitem{method}
\textit{Deletion compatibility and exhaustive extension search},
project documentation, \path{matroid1/docs/DELETION_COMPATIBILITY.md}.

\bibitem{status}
\textit{Completed rank-four search: status and computational resources},
project record, \path{matroid1/docs/STATUS.md}.

\bibitem{verification}
\textit{Rank-four, eleven-element verification record},
\path{matroid1/results/r4-n11/verification.json}.

\bibitem{drat}
N. Wetzler, M. J. H. Heule, and W. A. Hunt, Jr.,
\textit{DRAT-trim: Efficient checking and trimming using expressive clausal proofs},
Theory and Applications of Satisfiability Testing (SAT 2014),
Lecture Notes in Computer Science \textbf{8561}, 422--429.

\bibitem{diagnostic}
\textit{Diagnostic matroid instances and certificate validation},
project artifacts: \path{study/build_dataset.py},
\path{study/data/certification.json}, and
\path{study/output/lekang/validation/summary.json}.
The last record documents the 13 November 2026 recheck of all 24 cases.

% The supplied report cites classificationlogs and searchlogs above, but their
% corresponding log locations and bibliography entries have not been supplied.
% Add the real references here; do not substitute the separate diagnostic data.

\end{thebibliography}

\end{document}
